% ── AI Chat ───────────────────────────────────────────────────
\chapter{AI Chat}

        AI~chat je najinováčnejšou súčasťou systému. Umožňuje používateľom klásť
otázky v~prirodzenom jazyku nad obsahom nahraných dokumentov.
Systém odpovie na základe extrahovaného textu z~vybraných súborov, čím výrazne
urýchľuje orientáciu v~rozsiahlych podkladoch.

\section{Integrácia AI}

        Aplikácia podporuje dva AI~backendy, ktoré je možné konfigurovať podľa potrieb
organizácie:

        \textbf{Ollama (predvolený)} -- lokálny LLM server, ktorý beží na rovnakom
stroji ako aplikácia alebo na internom serveri organizácie. Predvolený model je
\texttt{gemma3:1b}, open-source model od spoločnosti Google. Výhodou je úplná
súkromnosť -- žiadne dáta neopúšťajú sieť organizácie. Nevýhodou je vyššia
hardvérová náročnosť servera.

        \textbf{Google Generative AI} -- cloudové API spoločnosti Google pre prípady,
keď lokálne prostriedky nepostačujú na beh väčšieho modelu. Pri použití tohto
backendu sa obsah otázky odosielala na externé servery.

  TBD

\section{Rozhranie chatu}

        TBD

\begin{figure}[H]
  \centering
  % \includegraphics[width=0.9\textwidth]{images/ai_chat.png}
  \fbox{\parbox{0.85\textwidth}{\centering\vspace{3cm}
    [TBD]
  \vspace{3cm}}}
  \caption{Rozhranie AI~chatu}
  \label{fig:ai_chat}
\end{figure}

        TBD

\begin{figure}[H]
  \centering
  % \includegraphics[width=0.9\textwidth]{images/ai_chat_odpoved.png}
  \fbox{\parbox{0.85\textwidth}{\centering\vspace{3cm}
    [TBD]
  \vspace{3cm}}}
  \caption{Príklad AI~odpovede nad obsahom dokumentu}
  \label{fig:ai_odpoved}
\end{figure}
